\paragraph{共振整函数的 Shannon 非可寻址性、局部 germ 传输与 Mellin 桥接猜想}
在共振整函数的 Cartwright 精确型、零点自相似递推与对数 Taylor 级数已经确定，而 fold-\(\pi\) 常数尾的奇指数 Bernoulli 层级亦已固定之后，还可把这两条主线继续压缩为下列七条接口结果。它们分别刻画整数栅格下的 Shannon 非可寻址性、零点计数的精确重整化、线性密度律、局部共振 germ 与 fold-\(\pi\) 奇层尾之间的 Bernoulli-free 传输、逐阶 Bernoulli 消去、有限共振射线族对 bulk 碰撞的不可压缩性，以及可能的 germ--Mellin 唯一桥接原理。

\begin{theorem}[共振整函数的 Shannon--Paley--Wiener 非可寻址性]\label{thm:xi-time-part9v-resonance-shannon-paley-wiener-obstruction}
记 \(\mathrm{PW}_\pi\) 为经典 Shannon--Paley--Wiener 空间。则
\[
\mathcal F_\varphi\notin \mathrm{PW}_\pi.
\]
更一般地，\(\mathcal F_\varphi\) 不属于任何以整数栅格 \(\ZZ\) 为稳定完全采样集的带限 \(L^2\) 整函数类。因而，整数共振样值
\[
u\longmapsto C_\varphi(u)=|\mathcal F_\varphi(u)|
\]
不是 Shannon 意义下可寻址的频谱对象。
\end{theorem}
\begin{proof}
由定理 \ref{thm:fold-pisot-bernoulli-convolution-representation} 的整数采样公式，
\[
|\mathcal F_\varphi(u)|=C_\varphi(u)
\qquad (u\in\ZZ).
\]
又由于 \(\mathcal F_\varphi\) 为偶函数且 \(\mathcal F_\varphi(0)=1\)，故
\[
\sum_{u\in\ZZ}|\mathcal F_\varphi(u)|^2
=
1+2\sum_{u=1}^{\infty}C_\varphi(u)^2.
\]
而定理 \ref{thm:fold-sigma-phi-diverges} 给出
\[
1+2\sum_{u=1}^{\infty}C_\varphi(u)^2=+\infty.
\]
因此
\[
\bigl(\mathcal F_\varphi(u)\bigr)_{u\in\ZZ}\notin \ell^2(\ZZ).
\]
然而，对任意 \(f\in \mathrm{PW}_\pi\)，经典 Shannon--Parseval 理论都给出
\(
(f(n))_{n\in\ZZ}\in \ell^2(\ZZ)
\)；
故 \(\mathcal F_\varphi\notin \mathrm{PW}_\pi\)。

更一般地，若 \(\mathcal B\) 是任意以 \(\ZZ\) 为稳定完全采样集的带限 \(L^2\) 整函数类，则样值算子
\[
f\longmapsto (f(n))_{n\in\ZZ}
\]
必连续取值于 \(\ell^2(\ZZ)\)。由于 \(\mathcal F_\varphi\) 的整数样值列不在 \(\ell^2(\ZZ)\)，故 \(\mathcal F_\varphi\notin \mathcal B\)。最后由 \(C_\varphi(u)=|\mathcal F_\varphi(u)|\) 即得结论。证毕。
\end{proof}

\begin{theorem}[共振零点计数的精确重整化方程]\label{thm:xi-time-part9v-resonance-zero-count-renormalization}
\leanverified{paper\_xi\_time\_part9v\_resonance\_zero\_count\_renormalization}
记
\[
N_\varphi(R):=\#\bigl(\mathcal Z(\mathcal F_\varphi)\cap[-R,R]\bigr)
\qquad (R\ge 0).
\]
则对每个 \(R\ge 0\) 都有严格恒等式
\[
N_\varphi(R)
=
N_\varphi\!\left(\frac{R}{\varphi}\right)
+
2\left\lfloor \frac{R}{\varphi^2}+\frac12\right\rfloor.
\]
等价地，
\[
N_\varphi(\varphi R)
=
N_\varphi(R)
+
2\left\lfloor \frac{R}{\varphi}+\frac12\right\rfloor.
\]
\end{theorem}
\begin{proof}
当 \(R=0\) 时结论显然成立。以下设 \(R>0\)。由定理 \ref{thm:fold-resonance-entire-lp} 的函数方程，
\[
\mathcal F_\varphi(z)
=
\cos\!\Bigl(\frac{\pi z}{\varphi^2}\Bigr)\,
\mathcal F_\varphi\!\Bigl(\frac{z}{\varphi}\Bigr).
\]
第一因子的零点恰为
\[
z=\left(k+\tfrac12\right)\varphi^2,
\qquad
k\in\ZZ,
\]
其落在区间 \([-R,R]\) 内的个数为
\[
2\left\lfloor \frac{R}{\varphi^2}+\frac12\right\rfloor.
\]

另一方面，定理 \ref{thm:fold-resonance-entire-lp} 还给出
\[
\mathcal Z(\mathcal F_\varphi)
=
\Bigl\{\left(k+\tfrac12\right)\varphi^n:\ k\in\ZZ,\ n\ge 2\Bigr\},
\]
故 \(\cos(\pi z/\varphi^2)\) 的零点对应 \(n=2\) 层，而 \(\mathcal F_\varphi(z/\varphi)\) 的零点对应 \(n\ge 3\) 层；不同层零点互不相交。于是 \([-R,R]\) 内的全部零点精确分解为两部分：一部分来自第一因子，另一部分来自第二因子。

对第二因子，\(\mathcal F_\varphi(z/\varphi)=0\) 当且仅当 \(z/\varphi\in \mathcal Z(\mathcal F_\varphi)\)。因此
\[
\#\Bigl(\mathcal Z\bigl(\mathcal F_\varphi(z/\varphi)\bigr)\cap[-R,R]\Bigr)
=
N_\varphi\!\left(\frac{R}{\varphi}\right).
\]
两部分相加即得
\[
N_\varphi(R)
=
N_\varphi\!\left(\frac{R}{\varphi}\right)
+
2\left\lfloor \frac{R}{\varphi^2}+\frac12\right\rfloor.
\]
再把 \(R\) 替换为 \(\varphi R\)，便得到等价的前向重整化式
\[
N_\varphi(\varphi R)
=
N_\varphi(R)
+
2\left\lfloor \frac{R}{\varphi}+\frac12\right\rfloor.
\]
证毕。
\end{proof}

\begin{corollary}[共振零点的线性密度律]\label{cor:xi-time-part9v-resonance-zero-density-linear}
\leanverified{paper\_xi\_time\_part9v\_resonance\_zero\_density\_linear}
沿用定理 \ref{thm:xi-time-part9v-resonance-zero-count-renormalization} 的记号，则
\[
N_\varphi(R)=2R+O(\log R)
\qquad (R\to\infty),
\]
并且
\[
\lim_{R\to\infty}\frac{N_\varphi(R)}{2R}=1.
\]
\end{corollary}
\begin{proof}
记
\[
E(R):=N_\varphi(R)-2R.
\]
由定理 \ref{thm:xi-time-part9v-resonance-zero-count-renormalization}，
\[
E(\varphi R)
=
E(R)+
2\left\lfloor \frac{R}{\varphi}+\frac12\right\rfloor
-
\frac{2R}{\varphi}.
\]
因此存在绝对常数 \(C>0\)，使得
\[
|E(\varphi R)-E(R)|\le C
\qquad (R>0).
\]
取 \(J:=\lfloor \log_\varphi R\rfloor\)，反复迭代直到 \(R\varphi^{-J}=O(1)\)，可得
\[
|E(R)|\le |E(R\varphi^{-J})|+CJ=O(\log R).
\]
于是
\[
N_\varphi(R)=2R+O(\log R).
\]
两边同除以 \(2R\) 并令 \(R\to\infty\)，便得到极限密度公式。证毕。
\end{proof}

\begin{theorem}[共振 germ 与 fold-\(\pi\) 奇层尾的 Bernoulli-free 传输律]\label{thm:xi-time-part9v-resonance-germ-foldpi-bernoulli-free-transfer}
在圆盘
\[
|z|<\frac{\varphi^2}{2}
\]
内写
\[
\log \mathcal F_\varphi(z)=-\sum_{r=1}^{\infty}b_r z^{2r}.
\]
再记 \(a_r\) 为定理 \ref{thm:fold-bin-gauge-constant-stirling-bernoulli-hierarchy} 中
\(
\log C_m-\log(2\pi)
\)
的第 \(r\) 个奇指数 Bernoulli 系数，即
\[
a_r=\frac{B_{2r}}{r(2r-1)}
\bigl(\varphi^{-1}+\varphi^{2r-3}\bigr),
\]
并定义
\[
\widetilde a_r:=(-1)^{r-1}a_r>0.
\]
则对每个 \(r\ge 1\)，都有
\[
b_r=
\frac{2^{2r}(2^{2r}-1)|B_{2r}|\pi^{2r}}{2r(2r)!}\,
\frac{\varphi^{-2r}}{\varphi^{2r}-1},
\]
并且满足严格比值恒等式
\[
\frac{\widetilde a_r}{b_r}
=
\frac{2(2r)!}{2^{2r}(2^{2r}-1)(2r-1)\pi^{2r}}\,
\bigl(\varphi^{-1}+\varphi^{2r-3}\bigr)\,
\varphi^{2r}(\varphi^{2r}-1).
\]
特别地，共振整函数在原点的局部 germ 与 fold-\(\pi\) 奇层尾共享同一 Bernoulli 骨架；二者之差完全由 \(\varphi\)-几何因子与 \(\pi^{-2r}\) 尺度转移给出。
\end{theorem}
\begin{proof}
由定理 \ref{thm:xi-time-part71-resonance-entire-log-taylor-zero-moments}，
\[
\log \mathcal F_\varphi(z)
=
-\sum_{r=1}^{\infty}\frac{\sigma_r}{r}\,z^{2r},
\qquad
\sigma_r=
(2^{2r}-1)\zeta(2r)\frac{\varphi^{-2r}}{\varphi^{2r}-1}.
\]
故 \(b_r=\sigma_r/r\)。再用 Euler 公式
\[
\zeta(2r)=\frac{2^{2r-1}|B_{2r}|\pi^{2r}}{(2r)!},
\]
便得到
\[
b_r=
\frac{2^{2r}(2^{2r}-1)|B_{2r}|\pi^{2r}}{2r(2r)!}\,
\frac{\varphi^{-2r}}{\varphi^{2r}-1}.
\]
另一方面，定理 \ref{thm:fold-bin-gauge-constant-stirling-bernoulli-hierarchy} 给出
\[
a_r=\frac{B_{2r}}{r(2r-1)}
\bigl(\varphi^{-1}+\varphi^{2r-3}\bigr).
\]
由于
\[
(-1)^{r-1}B_{2r}=|B_{2r}|,
\]
故
\[
\widetilde a_r=
\frac{|B_{2r}|}{r(2r-1)}
\bigl(\varphi^{-1}+\varphi^{2r-3}\bigr).
\]
将 \(\widetilde a_r\) 与 \(b_r\) 直接相除，Bernoulli 因子完全消去，即得所示比值公式。证毕。
\end{proof}

\begin{corollary}[Bernoulli 消去原理]\label{cor:xi-time-part9v-bernoulli-elimination-principle}
沿用定理 \ref{thm:xi-time-part9v-resonance-germ-foldpi-bernoulli-free-transfer} 的记号，对每个 \(r\ge 1\) 都有
\[
\frac{\widetilde a_r}{-[z^{2r}]\log \mathcal F_\varphi(z)}
\in
\QQ(\varphi,\pi).
\]
换言之，当 fold-\(\pi\) 奇层系数与共振整函数的局部 Maclaurin 系数逐阶相除时，全部 Bernoulli 数与全部偶值 \(\zeta(2r)\) 信息同时消失；残余完全退化为代数--几何尺度因子。
\end{corollary}
\begin{proof}
由定义
\[
-[z^{2r}]\log \mathcal F_\varphi(z)=b_r.
\]
再由定理 \ref{thm:xi-time-part9v-resonance-germ-foldpi-bernoulli-free-transfer}，
\(
\widetilde a_r/b_r
\)
已显式写成仅含 \(r,\varphi,\pi\) 的表达式，故结论成立。证毕。
\end{proof}

\begin{theorem}[有限共振射线族对 bulk 碰撞的不可压缩性]\label{thm:xi-time-part9v-finite-ray-family-bulk-collision-incompressibility}
设 \(U_m\subset\NN\) 为一列整数集合，并满足
\[
\sup_{m\ge 1}|U_m|<\infty.
\]
定义由该射线族承载的截断共振能量
\[
E_m(U_m):=1+2\sum_{u\in U_m}C_\varphi(u)^2.
\]
则有
\[
F_{m+2}\,\mathsf{Col}_m-E_m(U_m)\longrightarrow +\infty
\qquad (m\to\infty).
\]
特别地，任何只保留一致有界条整数共振射线的方案，都不可能在主尺度上饱和 bulk 碰撞。
\end{theorem}
\begin{proof}
记
\[
K:=\sup_{m\ge 1}|U_m|<\infty.
\]
由于 \(0\le C_\varphi(u)\le 1\)，对一切 \(m\) 都有
\[
E_m(U_m)\le 1+2K.
\]
另一方面，推论 \ref{cor:fold-collision-scale-diverges} 给出
\[
F_{m+2}\,\mathsf{Col}_m\longrightarrow +\infty.
\]
于是
\[
F_{m+2}\,\mathsf{Col}_m-E_m(U_m)
\ge
F_{m+2}\,\mathsf{Col}_m-(1+2K)
\longrightarrow +\infty.
\]
证毕。
\end{proof}

\begin{conjecture}[共振 germ 与 fold-\(\pi\) 奇层尾的 Mellin 唯一桥接原理]\label{conj:xi-time-part9v-resonance-germ-mellin-unique-bridge}
在本文框架内，若存在一条严格、无外参且可递阶延拓的桥接，把共振侧与 fold-\(\pi\) 侧直接联系起来，则该桥接不可能来自
\begin{enumerate}
  \item 整数栅格上的 Shannon 型采样；
  \item 任意基数一致有界的整数共振射线族；
  \item 任意把共振样值闭包压入 \(\ell^2\)-可控离散频谱的方案。
\end{enumerate}
与现有结果相容的候选机制应当是：先取
\[
-\log \mathcal F_\varphi(z)=\sum_{r=1}^{\infty}b_r z^{2r},
\]
再施加某个仅依赖于 \(\varphi\) 的 Mellin 型重整化算子 \(\mathfrak M_\varphi\)，使
\[
\mathfrak M_\varphi\bigl((b_r)_{r\ge 1}\bigr)
=
(a_r)_{r\ge 1}.
\]
\end{conjecture}
\begin{proof}[证据]
定理 \ref{thm:xi-time-part9v-resonance-shannon-paley-wiener-obstruction} 排除了 Shannon--Paley--Wiener 型桥接；定理 \ref{thm:xi-time-part9v-finite-ray-family-bulk-collision-incompressibility} 排除了基数一致有界的共振射线压缩；而定理 \ref{thm:xi-time-part9v-resonance-germ-foldpi-bernoulli-free-transfer} 与推论 \ref{cor:xi-time-part9v-bernoulli-elimination-principle} 又表明，在局部系数层面已经存在逐阶、显式且 Bernoulli 因子完全消失的传输比。因而，若严格桥接存在，则它不应发生在整数样值层，而应发生在局部 germ 经 Mellin 型重整化后的系数层。证毕。
\end{proof}

\noindent
上述结果表明，共振整函数 \(\mathcal F_\varphi\) 一方面处在 Laguerre--P\'olya/Cartwright 的临界边界上，另一方面却在整数栅格上失去 Shannon 可寻址性，并在 bulk 碰撞尺度上拒绝任何基数一致有界的共振射线压缩。与此同时，fold-\(\pi\) 奇层尾的系数与 \(-\log \mathcal F_\varphi\) 的局部系数又共享同一 Bernoulli 骨架，其差异完全退化为 \(\varphi\)-几何因子与 \(\pi^{-2r}\) 型尺度转移。因而，两侧之间真正可能的解析接口不再是整数样值的离散封装，而是局部 germ 经 Mellin 型重整化后的系数传输。

\endinput
